% SCP10, Lemma 6.2 and Observations 6.5--6.6, source lines 1704--1716 and 1825--1920.
% Bonds are unnormalized; one internal cycle contributes one factor of |G|.

\begin{theorem}[G-isometry under coordinate permutations]%
    \label{thm:peps_g_isometry_coordinate_equiv}
    \lean{TNLean.PEPS.IsGIsometric.of_coordinateEquiv}
    \leanok
    Let $T:\mathbb C^I\to\mathbb C^J$ be $G$-isometric for a
    representation $\rho$. Let $e:I'\to I$ and $f:J'\to J$ be bijections,
    and let $E:\mathbb C^{I'}\to\mathbb C^I$ and
    $F:\mathbb C^{J'}\to\mathbb C^J$ be their coordinate permutation
    maps. Suppose $E\rho'(g)=\rho(g)E$ for every $g$ and $FT'=TE$.
    Then $T'$ is $G$-isometric for $\rho'$, with the same positive
    isometry factor as $T$.
\end{theorem}

\begin{proof}\leanok
    The map $E$ carries invariant vectors to invariant vectors.
    The identity $FT'=TE$ transfers invariance and injectivity on the
    invariant subspace. Since $E$ and $F$ permute orthonormal bases,
    $\langle T'x,T'y\rangle=\langle TEx,TEy\rangle
        =c\langle Ex,Ey\rangle=c\langle x,y\rangle$ for invariant $x,y$.
\end{proof}

\begin{theorem}[Translation constraints around the internal square]%
    \label{thm:peps_two_by_two_translation_compatibility}
    \lean{TNLean.PEPS.twoByTwoSiteLegs_smul,
        TNLean.PEPS.twoByTwoSiteLegs_compatible_iff,
        TNLean.PEPS.sum_twoByTwoSiteLegs_compatible}
    \leanok
    \uses{def:peps_two_by_two_tensor}
    Let $G$ be a group. Write $\ell_i(\alpha,x)\in G^4$ for the
    four fine-leg configurations of Definition~\ref{def:peps_two_by_two_tensor},
    with corners indexed clockwise from the upper right. For paired boundary
    configurations $\alpha,\beta\in(G\times G)^4$ and internal labels
    $x,y,q\in G^4$, the equations
    $\ell_i(\alpha,x)=q_i\ell_i(\beta,y)$ at every corner hold precisely
    when $q_i=q_0$ at every corner, $\alpha=q_0\beta$, and $x=q_0y$.
    If $G$ is finite, then
    \begin{align}
        \sum_{x,y,q\in G^4}
        \prod_{i=0}^3\delta_{\ell_i(\alpha,x),q_i\ell_i(\beta,y)}
        =|G|^4\sum_{g\in G}\delta_{\alpha,g\beta}.
    \end{align}
\end{theorem}

\begin{proof}\leanok
    The internal bond shared by corners $0$ and $1$ gives
    $q_0y_1=x_1=q_1y_1$, hence $q_0=q_1$ by right cancellation.
    The bonds shared by corners $1,2$ and $2,3$ similarly give
    $q_1=q_2=q_3$. With this common translation, the sixteen fine-leg
    entries determine exactly the eight boundary labels and four internal
    labels. The converse follows by simultaneous translation of every label.
    For each $y\in G^4$ and each $g$ with $\alpha=g\beta$, the unique
    remaining assignment is $x=gy$ and $q_i=g$. There are $|G|^4$ choices
    of $y$.
\end{proof}

\begin{theorem}[Gram operator of the contracted four-site tensor]%
    \label{thm:peps_two_by_two_gram}
    \lean{TNLean.PEPS.twoByTwoTensor_gram_expansion,
        TNLean.PEPS.exists_twoByTwoTensor_gram}
    \leanok
    \uses{def:peps_two_by_two_tensor,thm:peps_regular_site_gram,
        thm:peps_two_by_two_translation_compatibility}
    Let $G$ and the physical alphabet $P$ be finite. Let the four tensors
    $a_i:G^4\times P\to\mathbb C$ be $G$-isometric for simultaneous
    regular translation. Denote their actual four-internal-bond contraction
    by $B(\alpha;\sigma)$, with $\alpha\in(G\times G)^4$ and
    $\sigma\in P^4$. There exist positive real numbers $c_i$ such that
    \begin{align}
        \sum_{\sigma\in P^4}
        \overline{B(\alpha;\sigma)}B(\beta;\sigma)
        =\left(\prod_{i=0}^3c_i\right)
        \sum_{g\in G}\delta_{\alpha,g\beta}.
        \label{eq:peps_two_by_two_gram}
    \end{align}
    Thus $B^*B=c_B P_8$, where $P_8$ averages simultaneous regular
    translation of the eight boundary bonds and
    $c_B=|G|\prod_i c_i>0$.
\end{theorem}

\begin{proof}\leanok
    The local isometries give
    $\sum_s\overline{a_i(\eta;s)}a_i(\theta;s)
        =(c_i/|G|)\sum_g\delta_{\eta,g\theta}$.
    Expanding both internal contractions and summing the four independent
    physical indices gives
    \begin{align}
        (B^*B)_{\alpha,\beta}
         & =\sum_{x,y\in G^4}\prod_{i=0}^3
        \sum_{s\in P}
        \overline{a_i(\ell_i(\alpha,x);s)}
        a_i(\ell_i(\beta,y);s)             \\
         & =\frac{\prod_i c_i}{|G|^4}
        \sum_{x,y,q\in G^4}
        \prod_{i=0}^3
        \delta_{\ell_i(\alpha,x),q_i\ell_i(\beta,y)}.
    \end{align}
    The translation-counting theorem cancels the factor $|G|^{-4}$ and
    gives (\ref{eq:peps_two_by_two_gram}). Since
    $(P_8)_{\alpha,\beta}=|G|^{-1}\sum_g\delta_{\alpha,g\beta}$,
    the block factor is $|G|\prod_i c_i$. The extra factor $|G|$ is
    contributed by the internal cycle with unnormalized bonds. This retains
    the scalar omitted from the diagrams in
    \cite[Lemma~6.2 and Observations~6.5--6.6]{Schuch2010PEPS}.
\end{proof}

\begin{theorem}[G-isometry of the explicit two-by-two contraction]%
    \label{thm:peps_two_by_two_isometry}
    \lean{TNLean.PEPS.twoByTwoTensor_translation,
        TNLean.PEPS.twoByTwoSeparatedTensor,
        TNLean.PEPS.isGIsometric_twoByTwoSeparatedTensor}
    \leanok
    \uses{thm:peps_two_by_two_gram,def:peps_two_by_two_tensor}
    Under the hypotheses of Theorem~\ref{thm:peps_two_by_two_gram},
    the literal contracted map $B:\mathbb C^{G^8}\to\mathbb C^{P^4}$
    is $G$-isometric for simultaneous regular translation of the eight
    boundary labels. The eight separate labels may equivalently be ordered
    as four paired legs. All four physical registers are retained, and
    surjectivity onto their ambient tensor product is not required.
\end{theorem}

\begin{proof}\leanok
    Translate all four internal summation variables by $g$.
    Local virtual invariance then gives $B(g\alpha;\sigma)=B(\alpha;\sigma)$.
    The Gram identity is $B^*B=c_BP_8$ with $c_B>0$, so
    $c_B^{-1}B^*$ is a left inverse on the invariant boundary subspace.
    For invariant $x,y$, it also gives
    $\langle Bx,By\rangle=c_B\langle x,y\rangle$.
    Separating each paired leg merely reindexes these formulas.
\end{proof}

\begin{theorem}[Nonvanishing of the original fine and coarse states]%
    \label{thm:peps_two_by_two_nonzero}
    \lean{TNLean.PEPS.IsGIsometric.torusBondNetwork_one_ne_zero,
        TNLean.PEPS.IsGIsometric.twoByTwoFineState_ne_zero,
        TNLean.PEPS.IsGIsometric.graphBondNetwork_torusIncidentFamily_ne_zero}
    \leanok
    \uses{thm:peps_g_isometry_coordinate_equiv,
        thm:peps_regular_torus_closure_independence}
    Let $a:G^4\times P\to\mathbb C$ be a regular $G$-isometric tensor,
    with $G$ and $P$ finite. Its homogeneous identity-closure state is
    nonzero on every torus with positive periods. In particular, the fine
    state on a $2w\times2h$ torus is nonzero. If $w,h\geq3$, the actual
    coarse graph state obtained by placing the same original tensor $a$
    at every coarse vertex is also nonzero. Neither conclusion requires
    surjectivity onto the ambient physical space.
\end{theorem}

\begin{proof}\leanok
    Reindex the four regular labels as an ordered quadruple. The coordinate
    bijection intertwines simultaneous translation, so it preserves
    $G$-isometry. The closure-sector nonvanishing theorem applies to this
    same tensor with identity closures. Its squared norm is positive because
    the common centralizer contains the identity. Identity closures put
    the identity matrix on every seam, giving the asserted native bond
    contraction. For $w,h\geq3$, its coefficient equals the incident-graph
    contraction coefficient under the four-leg enumeration.
\end{proof}
